\documentclass[11pt,a4paper]{article}
\usepackage[T1]{fontenc}
\usepackage[utf8]{inputenc}
\usepackage{lmodern,amsmath,amssymb,amsthm,booktabs,microtype}
\usepackage[margin=28mm]{geometry}
\usepackage[colorlinks=true,linkcolor=black,citecolor=blue,urlcolor=blue]{hyperref}
\hypersetup{pdftitle={Exact Independent-Set Counts in Motzkin Graphs through Order 144},pdfauthor={Carlo Corti}}
\newtheorem{theorem}{Theorem}[section]
\newtheorem{proposition}{Proposition}[section]
\newtheorem{corollary}{Corollary}[section]
\newcommand{\MG}{\mathrm{MG}}
\newcommand{\IS}{\mathcal{I}}
\title{Exact Independent-Set Counts in Motzkin Graphs\\through Order 144}
\author{Carlo Corti\\Independent researcher\\\texttt{carlo.corti@outlook.com}}
\date{September 2026}
\begin{document}
\maketitle

\begin{abstract}
We enumerate the independent vertex subsets, including the empty set, of the
Motzkin Riordan graph $\MG_n$ for every $1\leq n\leq144$.  A C17 count by
vertex deletion agrees at each of these 144 indices with a separately
implemented Python count based on independent odd-vertex subsets and their
available even vertices.  The programs also construct adjacency by different
arithmetic: truncated Motzkin-series powers over $\mathbb F_2$ and exact
integer coefficients from Lagrange inversion, respectively.  Relative to
the twelve-term table of Cheon and coauthors and the saved OEIS A339592
entry, the locally validated extension comprises indices 13--144, beginning
with $a(13)=104$ and ending with
$a(144)=4839167564472933081443$.  We prove the counting identities and a
finite 128-bit arithmetic bound, apply a Pell-number upper bound to the
Motzkin graph, identify shared assumptions and distinguish
the computed range from the uncomputed order 145.  The extended b-file,
programs and supporting evidence are publicly available with the project
repository and its Zenodo archive.
\end{abstract}
\noindent\textbf{2020 Mathematics Subject Classification:} Primary 05C30;
secondary 05C69, 05A15.\\
\textbf{Key words and phrases:} Motzkin graph, Riordan graph, independent set,
exact enumeration, integer sequence, computational validation.

\section{Introduction}\label{sec:intro}
The Motzkin graph $\MG_n$ is a labelled Riordan graph on $[n]=\{1,\ldots,n\}$.
Let $a(n)$ count all its independent vertex subsets, including the empty set.
Cheon, Jung, Kang, Kim, Kim, Kitaev and Mojallal gave the values
\[
2,3,4,7,9,13,17,26,29,48,55,95
\]
for $1\leq n\leq12$ in their Table~1~\cite{counting}.  Their concluding
section asked for general bounds or exact enumeration for the Motzkin
family.  N.~J.~A. Sloane
submitted the corresponding OEIS entry A339592 on December 28, 2020; the
local snapshot consulted in September 2026 has precisely those twelve terms,
offset $1,1$, and the keyword
\texttt{more}~\cite{oeis}.  The underlying Riordan-graph construction and
structural framework predate this computation~\cite{structure}.

This article records a finite extension and its local validation.
The C17 campaign enumerated each order from 1 through 144.  An independently
implemented Python algorithm compared every indexed value over the same
range.  The local b-file was then updated from 1--12 to 1--144; the saved OEIS
entry was not edited.  No order beyond 144 is admitted from this campaign.
The main distinction from the Pascal and Catalan companion graphs is that
$\MG_n$ is not io-decomposable~\cite{structure}; their odd-subgraph
self-similarity and universal-vertex consequences do not provide Motzkin
values.  The even-vertex null subgraph is a previously established Bell-type
property~\cite{structure}.  Here it yields a different,
weighted enumeration, while the production program uses general vertex
deletion.  The exact coefficient formulas and an arithmetic bound make both
routes explicit.

\section{Definition and structural identities}\label{sec:def}
Let $M(z)=\sum_{r\geq0}M_rz^r$ be the unique formal power series satisfying
\begin{equation}\label{eq:motzkin}
 M(z)=1+zM(z)+z^2M(z)^2,\qquad M_0=1.
\end{equation}
The Motzkin graph is the Bell-type Riordan graph
$\MG_n=G_n(M,zM)$~\cite{structure,counting}.  It is simple and undirected,
with zero diagonal and, for $1\leq j<i\leq n$,
\begin{equation}\label{eq:edge}
 A_{ij}=A_{ji}=\bigl([z^{i-j-1}]M(z)^j\bigr)\pmod{2}.
\end{equation}
This follows from the Riordan rule
$A_{ij}=[z^{i-2}]g(z)f(z)^{j-1}\pmod{2}$ after setting
$g=M$ and $f=zM$.  The exponent and labels are both one-based: the
coefficient degree is $i-j-1$, not $i-j$.  The first $m$ vertices of
$\MG_n$ induce $\MG_m$ for $m\leq n$.  We count subsets of this labelled
graph, without quotienting by automorphisms.  Put
$\IS(G)=\{S\subseteq V(G):S\text{ is independent}\}$ and $i(G)=|\IS(G)|$;
then $a(n)=i(\MG_n)$ and $i(\varnothing)=1$.  For example,
$\MG_3=K_3$, so $a(3)=4$.

\begin{proposition}\label{prop:parity}
Every pair of consecutive vertices of $\MG_n$ is adjacent, and the
even-labelled vertices $E_n=\{2,4,\ldots,2\lfloor n/2\rfloor\}$ form an
independent set.  Consequently, with $F_0=0,F_1=1$,
\begin{equation}\label{eq:bounds}
 2^{\lfloor n/2\rfloor}\leq a(n)\leq F_{n+2}.
\end{equation}
\end{proposition}
The even-vertex assertion is the Bell-type o-decomposition of
Cheon et al.~\cite[Theorem~4.3]{structure}; their decomposition also
describes the odd-induced graph by a different Riordan pair.  The path
bound appears in~\cite[Theorem~3.2 and Corollary~3.3]{counting}, where
the authors use $F(0)=F(1)=1$, so their $F(n+1)$ is our $F_{n+2}$.
We include the short special-case proof to fix the parity and index
conventions used in the computations.
\begin{proof}
For $i=j+1$, the coefficient in~\eqref{eq:edge} is the constant term
of $M^j$, namely 1.  If $i$ and $j$ are even, $d=i-j-1$ is odd and,
over $\mathbb F_2$, $M(z)^j=(M(z)^{j/2})^2$ has no odd-degree term:
squaring a binary power series doubles every exponent.
Every subset of $E_n$ is therefore independent, giving the lower bound.
Consecutive edges force every independent-set indicator word to avoid
consecutive ones; the number of such length-$n$ words is $F_{n+2}$.
\end{proof}

Cheon et al. used an auxiliary graph $\Delta_n$ to bound the counts of
io-decomposable Bell-type graphs~\cite{counting}.  The same edge check
applies to $\MG_n$ without assuming io-decomposability.  Its resulting
bound is at least as sharp as the path bound and is strict for $n\geq3$.
Let $P_0=0$, $P_1=1$ and
$P_{k+2}=2P_{k+1}+P_k$ be the Pell numbers.  Define
\[
 \delta_n=
 \begin{cases}
  2P_{(n+1)/2},&n\text{ odd},\\
  P_{n/2}+P_{n/2+1},&n\text{ even}.
 \end{cases}
\]
\begin{proposition}\label{prop:pell}
For every $n\geq1$,
\[
 2^{\lfloor n/2\rfloor}\leq a(n)\leq\delta_n\leq F_{n+2}.
\]
\end{proposition}
\begin{proof}
Let $\Delta_n$ have all edges $\{j,j+1\}$ and all
$\{j,j+2\}$ with odd $j$ and $j+2\leq n$.  The path edges occur by
Proposition~\ref{prop:parity}.  Since $M_1=1$, for odd $j$ we have
$[z]M(z)^j=jM_1\equiv1\pmod{2}$; hence the odd-step edges also occur
in $\MG_n$.  Thus $\Delta_n$ is a spanning subgraph of $\MG_n$.
Cheon et al.~\cite[Lemma~3.4]{counting} count the independent sets of
$\Delta_n$ as $\delta_n$ with the displayed Pell normalization.
Deleting edges cannot decrease the independent-set count, so
$a(n)\leq i(\Delta_n)=\delta_n$.  The remaining bounds follow from
Proposition~\ref{prop:parity} and the fact that $\Delta_n$ contains
the path on $[n]$.  For $n\geq3$, the path-independent set
$\{1,3\}$ is excluded by $\Delta_n$, so $\delta_n<F_{n+2}$.
\end{proof}

The inclusion is not peculiar to Motzkin graphs: for any proper Bell-type
$G_n(g,zg)$ with odd $g_1$, its path edges and
$[z]g(z)^j=jg_1\equiv1\pmod{2}$ for odd $j$ contain $\Delta_n$.
Here we apply that edge argument only to $\MG_n$; the stronger
io-decomposable bounds of~\cite{counting} are not invoked.

Write $O_n=[n]\setminus E_n$ and let $N(T)$ be the union of the open
neighbourhoods of vertices in $T$.  The structural fact in
Proposition~\ref{prop:parity} gives the counting identity actually used by
the second program.
\begin{proposition}\label{prop:weighted}
For every $n\geq1$,
\begin{equation}\label{eq:weighted}
 a(n)=\sum_{T\in\IS(\MG_n[O_n])}
 2^{|E_n\setminus N(T)|}.
\end{equation}
\end{proposition}
\begin{proof}
An independent set $S$ has a unique odd part $T=S\cap O_n$, which must be
independent in the odd-induced graph.  Once $T$ is fixed, its even part can
be any subset of $E_n\setminus N(T)$, because no two even vertices are
adjacent.  The displayed summands count disjoint classes and include
$T=\varnothing$ and the empty set.
\end{proof}

The odd-induced graph is not assumed to be another Motzkin graph.
In $\MG_5[\{1,3,5\}]$, vertices 1 and 3 are adjacent since $M_1=1$,
and 3 and 5 are adjacent since $[z]M(z)^3=3$; vertices 1 and 5 are
nonadjacent since $M_3=4$ is even.  This induced graph is the path $P_3$,
whereas $\MG_3=K_3$, so the two graphs are not isomorphic.  This
directly illustrates the non-io-decomposability established
in~\cite[Example~4.8]{structure}.  O-decomposability supplies the
even-null fact but does not imply io-decomposability.  No
universal-vertex formula is used below.

\section{Two exact graph and count constructions}\label{sec:alg}
\subsection{C17 construction and deletion count}
All C adjacency arithmetic is in $\mathbb F_2$.  Writing
$b_r=M_r\pmod{2}$, Equation~\eqref{eq:motzkin} and Frobenius give
\begin{equation}\label{eq:parityrec}
 b_0=1,\qquad
 b_r=b_{r-1}+\mathbf1_{2\mid r,\ r\geq2}\,b_{(r-2)/2}
 \quad(r\geq1),
\end{equation}
where addition is modulo 2.  The C implementation builds the truncated
powers $M^j$ over $\mathbb F_2$ by successive convolution with the $b_r$ and applies
\eqref{eq:edge}.  The graph is held as symmetric adjacency bit masks.
Here Frobenius means that, for $B(z)=\sum c_rz^r$ over $\mathbb F_2$,
$B(z)^2=\sum c_rz^{2r}$.

For any induced vertex set $S$ and $v\in S$, let
$I(S)=i(\MG_n[S])$ denote its
independent-set count and $N_S[v]$ its closed neighbourhood within $S$.
The C kernel uses
\begin{equation}\label{eq:deletion}
 I(\varnothing)=1,\qquad
 I(S)=I(S\setminus\{v\})+I(S\setminus N_S[v]).
\end{equation}
It chooses a vertex of maximum degree in the current induced graph,
memoizes counts under the full mask, and handles an edgeless state of size
$k$ as $2^k$.  The two terms partition independent sets according to
whether $v$ is excluded or included.  Each recursive child has fewer
vertices, proving termination and the exactness of the recurrence; cache
collisions replace entries but do not change the mask comparison or the
returned count.  Counts use checked unsigned 128-bit additions and three
64-bit mask limbs.  The campaign source was compiled by its C17 Makefile
with GCC 13 at \texttt{-O3 -std=c17} and the warning flags
\texttt{-Wall -Wextra -Wpedantic -Wshadow -Wconversion}
and \texttt{-Wdouble-promotion -Wformat=2}.  The Makefile supplies
\texttt{SOURCE\_SHA} and \texttt{MAKE\_SHA} from the source and recipe
bytes; those macros are required for the source to compile.  The
executable embeds the resulting identities in its run metadata.
Its \texttt{run}, \texttt{resume} and
\texttt{verify} modes recover or inspect complete indexed terms, not partial
recursion states.  The \texttt{verify} mode checks state integrity and
executable identity; it does not recount the graph.

\begin{proposition}\label{prop:capacity}
For $1\leq n\leq192$, the true value $a(n)$ is smaller than $2^{128}$.
The same bound applies to every induced state counted by the C kernel.
\end{proposition}
\begin{proof}
Consider six consecutive labelled vertices.  Their internal edge
coefficients have degree at most 4.  Shifting both labels by eight
multiplies the relevant power of $M$ by $M^8=M(z^8)$ over
$\mathbb F_2$, which is 1 through degree 7.  Thus the induced
six-vertex graph depends only on its first label modulo 8.
Direct application of~\eqref{eq:parityrec} and~\eqref{eq:edge},
followed by exhaustive testing of the $2^6$ subsets, gives
\[
\begin{array}{c|rrrrrrrr}
\text{first label }s&1&2&3&4&5&6&7&8\\ \hline
i(\MG_{s+5}[\{s,\ldots,s+5\}])&13&14&13&14&15&12&14&15
\end{array}
\]
The C self-test checks these same eight counts; the finite subset
calculation is separate from the period-eight proof.  Partition $[192]$ into
32 consecutive six-vertex blocks.  Restriction of an independent set
to the blocks is injective, and each block has at most 15 choices, so
$a(192)\leq15^{32}=(225)^{16}<(256)^{16}=2^{128}$.
As a separate, sharper bound,
Proposition~\ref{prop:pell} gives
$a(192)\leq\delta_{192}
=6733044458057842709277507685523012161<2^{128}$.
Since $\MG_n$ is an induced prefix,
$a(n)\leq a(192)$ for $n\leq192$.  Every independent set of an induced
state is also an independent set of $\MG_n$, giving the final assertion.
\end{proof}

This is a numeric representation bound, not a runtime guarantee.  The
three-mask-limb design accepts orders through 192, but neither it nor the
bound supplies a count for 145--192.  The implementation rejects arithmetic
overflow rather than silently wrapping.

\subsection{Python construction and weighted count}
The Python checker reconstructs the adjacency independently of the C
parity recurrence.  Put $w=zM(z)$; Equation~\eqref{eq:motzkin} becomes
$w=z(1+w+w^2)$.  Lagrange inversion~\cite{gessel} gives, for
$j\geq1$ and $d\geq0$,
\begin{equation}\label{eq:lagrange}
 [z^d]M(z)^j
 =\frac{j}{d+j}[u^d](1+u+u^2)^{d+j}
 =\frac{j}{d+j}\sum_{k=0}^{\lfloor d/2\rfloor}
 \binom{d+j}{k}\binom{d+j-k}{d-2k}.
\end{equation}
For $d=0$ both sides equal 1.  For $d>0$, apply Lagrange inversion to
$w^j=z^jM^j$ at degree $d+j$.  Python evaluates the integer sum exactly,
checks divisibility by $d+j$, and reduces the result modulo 2 only when
forming an edge.  It then verifies the even-even null condition and
enumerates independent odd subsets recursively.  At each completed odd
subset it adds $2^{|E_n\setminus N(T)|}$, as in~\eqref{eq:weighted}.
The odd recursion either excludes its next vertex or includes it only
when no selected odd vertex is adjacent; it therefore visits each valid
$T$ once and terminates after $\lceil n/2\rceil$ decisions.

The programs consequently differ both in coefficient arithmetic and in
the counting partition.  They still share Equation~\eqref{eq:edge}, the
one-based vertex convention, and the interpretation of an independent set.
Agreement cannot detect an identical mistake in those shared premises.
The original 1--12 table anchors the graph convention externally.
Before the official campaign, a disposable direct enumeration tested all
$2^n$ subsets at each $1\leq n\leq13$, reproducing that prefix and
obtaining 104 at order 13.  Its local diagnostic report records this
finite check, which shares its coefficient-derived graph with the
contemporaneous weighted check and was not an admitted campaign result.

\section{Local campaign, validation and provenance}\label{sec:results}
The frozen C17 source, build recipe and executable produced one native run
with target 144.  Its indexed snapshot has exactly one row for every
$n=1,\ldots,144$; the native log has one start and one commit per order and a
completed 144-row terminal record.  The Python 3.12 standard-library checker
read that frozen snapshot and compared its own count with every indexed row
from 1 through 144.  Its separate state and log record 144 completed
comparisons and zero unchecked rows.  The twelve old values match both
Table~1 of~\cite{counting} and the saved OEIS entry~\cite{oeis}.
The two exact programs' complete indexed agreement is the basis for the
locally validated computational claim; the logs and checksums establish
coverage and file identity, rather than mathematical correctness by
themselves.  The original snapshot, checker state, Makefile and logs are
retained with the local project but have not yet been distributed in a
public archive.  Project source, code-review records and local result evidence
underlie this account~\cite{project}.

\begin{theorem}[Finite computational result]\label{thm:finite}
With the graph and indexing of~\eqref{eq:edge}, the C17 campaign and the
distinct Python algorithm give identical exact integers $a(n)$ at each
index $1\leq n\leq144$.  In particular, the 132 locally validated values
at indices $13\leq n\leq144$ extend the twelve-term reference prefix, and
\[
 a(13)=104,\qquad a(144)=4839167564472933081443.
\]
The first index without a result from this campaign is 145.
\end{theorem}
\begin{proof}[Evidence and computational justification]
Equations~\eqref{eq:parityrec} and~\eqref{eq:edge} specify the C graph;
\eqref{eq:deletion} proves that its terminating count enumerates every
independent set once.  Equation~\eqref{eq:lagrange} gives a separate exact
edge construction, and Proposition~\ref{prop:weighted} proves the Python
count.  The frozen native snapshot contains consecutive rows 1--144, and the
checker state records a matching row at each of those indices.  Both logs
close their stated ranges without an unchecked index.  The b-file comparison
also matches those rows by mathematical index, but the later b-file is not a
third independent count.  These locally audited records supply the
finite computational evidence, subject to the shared-definition and
implementation limits described above.  An integrity check alone would not
establish any of the large-order counts.
\end{proof}

\begin{table}[ht]
\centering
\begin{tabular}{r@{\qquad}r@{\qquad}l}
\toprule
$n$ & $a(n)$ & Status relative to the saved entry \\
\midrule
12 & 95 & Previous table and snapshot \\
13 & 104 & Locally enumerated and compared \\
64 & 4352673893 & Locally enumerated and compared \\
128 & 18471061891119082367 & Locally enumerated and compared \\
144 & 4839167564472933081443 & Locally enumerated and compared \\
\bottomrule
\end{tabular}
\caption{Selected indexed values; the complete local b-file contains 1--144.}
\label{tab:values}
\end{table}

The native run took about 13 seconds of wall time and reported a peak
resident set of 50,684 KiB; the Python run took about 266 seconds and
reported 18,612 KiB.  These measurements come from the original run logs.
The project's pre-run environment record lists an AMD Ryzen 9 7940HS Linux
host and Python 3.12.3, while the Makefile specifies GCC 13; the environment
record is context rather than an independent runtime trace.  Neither
implementation launches parallel workers.  These timings are not asymptotic
estimates.  The final two Python terms each took about 96 seconds in the
recorded cumulative run.  Earlier isolated diagnostics used a different
invocation context; the cause of the timing difference has not been
established and has no bearing on the exact-count comparison.  The observed
timings do not support a
projection to order 160 or 192.  The campaign did not execute those orders.

The chronology of local data matters.  The result-validation report was
written when the b-file still held 1--12.  A subsequent, separately recorded
update copied the complete validated native rows 1--144 into the local
b-file, preserving the former twelve-row file as historical evidence.  The
new b-file has 144 consecutive indexed rows and two terminal line-feed
bytes, representing a final blank line.
Neither this update nor the computational validation changed the saved OEIS
snapshot or constitutes an OEIS submission.

\section{Finite structure and limits}\label{sec:limits}
The induced-prefix relation yields one general cross-order consequence:
$a(n+1)>a(n)$ for every $n\geq1$, because each independent set on $[n]$
remains one on $[n+1]$ and the singleton $\{n+1\}$ adds a distinct set.
It also gives
$a(n+1)\leq a(n)+(a(n)-1)<2a(n)$: of the old independent
sets, at least $\{n\}$ cannot be extended by vertex $n+1$, because
consecutive vertices are adjacent.  The 144 indexed values satisfy both inequalities.
These statements are proved from the graph definition, not conjectured from
the finite table.

The even-null identity supplies the lower bound and weighted checker, but
does not turn the odd-induced graph into $\MG_{\lceil n/2\rceil}$.
The structural source explicitly excludes Motzkin graphs from the relevant
io-decomposable class~\cite{structure}.  The initial Pascal, Motzkin and
Catalan rows share a Riordan framework~\cite{counting}, but their defining
series and indexed counts differ.  Comparing the available finite data
did not yield a verified cross-family identity, general recurrence or
asymptotic formula for $a(n)$.  No such claim is made.  In particular, a
dyadic universal-vertex identity from another graph family cannot be used
to produce $a(145)$ here.  Resolving order 145 would require a new,
completed enumeration with corresponding validation; additional numeric
type capacity alone would not do so.

\section{Conclusion}
The local computation supplies 132 validated additions to the twelve-term
A339592 reference prefix and a complete indexed b-file through 144.
The endpoint value is $a(144)=4839167564472933081443$; order 145 remains
unresolved by this campaign.  The previously known even-vertex null
subgraph enabled a
second exact enumeration with coefficient arithmetic and counting logic
distinct from the production program.  This validates a substantial finite
range.  Applying the previously used auxiliary-graph method to $\MG_n$
yields the general upper bound $a(n)\leq\delta_n$; it does not give
an exact formula.  General Motzkin-graph enumeration remains open.

\section*{Data availability}
The C source and build recipe, Python checker, indexed native results,
checker state, original logs, pre-run environment record, small-order
diagnostic and validation reports, complete 1--144 b-file, and this
manuscript's TeX source and PDF are available in the public project repository
\url{https://github.com/carcorti/A339592}.  GitHub releases are archived
under the Zenodo concept DOI
\url{https://doi.org/10.5281/zenodo.23010698}.  This concept DOI identifies
the release family; each deposited version has its own record DOI.  The
included validation command audits the retained 1--144 results and recomputes
small orders, but does not repeat the full campaign.  The saved OEIS snapshot
still records 1--12, and this work supplies no exact value beyond 144.

\section*{AI use disclosure}
Large language models were used as research and drafting assistants during
the A339592 workflow, including OpenAI Codex and external review systems
whose local reports carry the interface or family labels ChatGPT, Claude,
DeepSeek, Gemini, GLM, Grok, MiniMax, Muse and Qwen.  These labels do not
independently establish the serving model, version or backend for every
review.  No AI system is an author.

\paragraph{Mathematical analysis.} AI assistance included literature
search, comparison of the Riordan-graph sources and scrutiny of draft
arguments.  The published Bell-type decomposition and auxiliary-graph
counts are credited to their authors.  The deductions for $\MG_n$,
definitions and finite bounds are stated with proofs or explicit
calculations rather than on the authority of a model.

\paragraph{Algorithm and code.} The project records identify Codex
assistance in the development and hardening of both the C program and the
Python checker; they do not establish reliable line-by-line tool authorship.
Reviewers supplied code observations and bounded diagnostic programs.
A reproduced persistence defect in the earlier C candidate led to the
reviewed C source used here.  Frozen source and build identities,
self-tests, original run records and the Python per-index comparison
support the computational claims.  The two programs are methodologically
distinct but share the graph definition, indexing and a project
development process; their agreement is not model-independent proof.

\paragraph{Manuscript.} AI assistance was used for drafting, source checks,
internal editorial review and the external review reports considered in
this revision.  Reviewer assertions and submitted diagnostic scripts were
weighed against primary sources and local evidence, not accepted as
execution records.  Final numerical claims rest on indexed artifacts and
validation rather than conversational assertions.
Carlo Corti bears full scientific responsibility for the mathematical
statements, algorithm, implementation, numerical results, validation
boundaries and conclusions.

\clearpage
\begin{thebibliography}{9}
\sloppy\raggedright
\bibitem{structure}
G.-S. Cheon, J.-H. Jung, S. Kitaev and S. A. Mojallal,
``Riordan graphs I: Structural properties,''
\emph{Linear Algebra and its Applications} \textbf{579} (2019), 89--135.
\url{https://doi.org/10.1016/j.laa.2019.05.033}.
Theorem numbering cited here follows the accessible arXiv version,
\url{https://arxiv.org/abs/1710.04604v2}.

\bibitem{counting}
G.-S. Cheon, J.-H. Jung, B. Kang, H. Kim, S.-R. Kim,
S. Kitaev and S. A. Mojallal,
``Counting independent sets in Riordan graphs,''
\emph{Discrete Mathematics} \textbf{343} (2020), no.~11, article 112043.
\url{https://doi.org/10.1016/j.disc.2020.112043}.
Theorem and lemma numbering cited here follows the accessible arXiv version,
\url{https://arxiv.org/abs/2006.16579v1}.

\bibitem{oeis}
N. J. A. Sloane, ``A339592: Number of independent sets in Motzkin graph
on $n$ nodes,'' \emph{The On-Line Encyclopedia of Integer Sequences},
submitted December 28, 2020; local snapshot consulted September 2026.
\url{https://oeis.org/A339592}.

\bibitem{gessel}
I. M. Gessel, ``Lagrange inversion,'' \emph{Journal of Combinatorial
Theory, Series A} \textbf{144} (2016), 212--249.
Journal DOI: \url{https://doi.org/10.1016/j.jcta.2016.06.018}.
Theorem 2.1.1 was checked in the accessible arXiv version,
\url{https://arxiv.org/html/1609.05988v1}.

\bibitem{project}
C. Corti, \emph{Exact Independent-Set Counts in Motzkin Graphs through
Order 144}, GitHub repository and Zenodo archive, 2026.
\url{https://github.com/carcorti/A339592};
\url{https://doi.org/10.5281/zenodo.23010698}.
The Zenodo link is the concept DOI for the project release family, not a
version-specific record DOI.
\end{thebibliography}
\end{document}
